\section{Received-Power-Aware Adaptive Carrier Phase Recovery}
\label{sec:method}

The adaptive receiver contains DA and NDA carrier-recovery branches. Its control statistics are formed from each $N_{\mathrm{DSP}}=256$-sample raw received window before branch-specific equalization or phase estimation. Once the two-stage rule selects a branch, the receiver executes only that branch and returns one 256-sample carrier-corrected complex sequence to the demodulator. Figure~\ref{fig:adaptive-cpr} shows this select-before-execute data flow.

\subsection{DA and NDA Carrier-Recovery Algorithms}

The DA branch uses one known pilot every four transmitted symbols. Let $n_\ell$ denote the index of the $\ell$th pilot in the current processing window, $p_\ell$ its known complex value, and $t_\ell=n_\ell T_s$ its sampling time. Removing the known pilot modulation and fitting a linear phase model gives
\begin{equation}
\begin{aligned}
\vartheta_\ell
&=\operatorname{unwrap}\!\left\{\angle\!\left(\frac{r_{n_\ell}}{p_\ell}\right)\right\},\\
\vartheta_\ell
&\simeq\phi_0+2\pi\Delta f\,t_\ell,\\
\widehat{\Delta f}_{\mathrm{DA}}
&=\frac{\sum_\ell(t_\ell-\bar t)(\vartheta_\ell-\bar\vartheta)}
{2\pi\sum_\ell(t_\ell-\bar t)^2},\\
\hat\phi_{0,\mathrm{DA}}
&=\bar\vartheta-2\pi\widehat{\Delta f}_{\mathrm{DA}}\bar t .
\end{aligned}
\label{eq:da-phase-time-ls}
\end{equation}
The 64 regularly spaced pilots in a 256-sample window jointly determine the frequency slope and phase intercept. Removing $\hat\phi_{0,\mathrm{DA}}+2\pi\widehat{\Delta f}_{\mathrm{DA}}kT_s$ yields a carrier-corrected complex sequence with the same length and timing as the input. The pilot-energy convention follows Eq.~\eqref{eq:pilot-energy-coordinate}.

The NDA branch removes the $(8,8)$-16APSK modulation periodicity through the $M_0$th-power operation, with $M_0=8$, following the corresponding M-APSK construction~\cite{du2025nda}. A frequency estimate is first obtained from the spectrum of the raised signal and removed from the received window. Carrier phase is then estimated from the argument of the averaged $M_0$th-power samples. For a sample set $\mathcal I$, the basic estimate is
\begin{equation}
\hat\phi_{\mathcal I}
=\frac{1}{M_0}\angle\!\left(
\frac{1}{|\mathcal I|}\sum_{k\in\mathcal I}r_k^{M_0}
\right),\qquad M_0=8.
\label{eq:nda-m0-phase}
\end{equation}
The NDA frequency stage first locates the eighth-power spectral component, divides its frequency by $M_0$, and removes the resulting phase ramp. Under Gamma--Gamma turbulence, $\mathcal I$ is the complete 256-sample window. In the additive white Gaussian noise configuration, eight 32-sample segments provide raised-domain phase estimates that are unwrapped and linearly interpolated. Both modes derotate the frequency-compensated samples and return a 256-sample complex window.

The eighth-power statistic is defined modulo $2\pi/8$. For BER evaluation only, each window tests the eight constellation rotations against the transmitted labels and retains the rotation with the lowest error count; the branch controller and NDA phase statistic do not consume those labels.

\subsection{Two-Stage Branch Control}

The controller processes the raw received power $q_k=|r_k|^2$ over the current 256-sample window. Its first statistic is the CV:
\begin{equation}
\begin{aligned}
\mathrm{CV}&=\frac{\operatorname{std}(q_k)}{\operatorname{mean}(q_k)},\\
\tau_{\mathrm{CV}}(\gamma_{\mathrm{data}}^{\mathrm{dB}})
&=1.10\left(0.74+0.12e^{-\gamma_{\mathrm{data}}^{\mathrm{dB}}/5}\right).
\end{aligned}
\label{eq:cv-decision-boundary}
\end{equation}
CV measures power dispersion relative to the window mean: a low value identifies a comparatively stable window, whereas a high value triggers the second-stage test for fading-induced power variation. Here, $\gamma_{\mathrm{data}}^{\mathrm{dB}}$ is the nominal data-symbol SNR. The reference curve was fitted to an offline AWGN characterization, and $1.10$ is the fixed margin. Windows below the boundary activate NDA; the remainder proceed to the effective-SNR test.

For the second stage, the linear nominal data-symbol SNR is $\bar\gamma=E_s/N_0=10^{\gamma_{\mathrm{data}}^{\mathrm{dB}}/10}$. The controller forms a blind power-based channel proxy, maps it to an effective SNR, and applies the complete branch rule:
\begin{equation}
\begin{aligned}
\hat h_{\mathrm{dsp}}
&=\max\!\left\{
\frac{1}{N_{\mathrm{DSP}}}\sum_{k=0}^{N_{\mathrm{DSP}}-1}|r_k|^2
-\frac{1}{2\bar\gamma},\ 10^{-6}
\right\},\\
\hat\gamma_{\mathrm{eff,dB}}
&=\gamma_{\mathrm{data}}^{\mathrm{dB}}+10\log_{10}\hat h_{\mathrm{dsp}},\\
\mathcal A(r)
&=\begin{cases}
\mathrm{NDA}, & \mathrm{CV}<\tau_{\mathrm{CV}},\\
\mathrm{DA}, & \mathrm{CV}\ge\tau_{\mathrm{CV}}
\ \text{and}\ \hat\gamma_{\mathrm{eff,dB}}<13~\mathrm{dB},\\
\mathrm{NDA}, & \mathrm{CV}\ge\tau_{\mathrm{CV}}
\ \text{and}\ \hat\gamma_{\mathrm{eff,dB}}\ge13~\mathrm{dB}.
\end{cases}
\end{aligned}
\label{eq:two-stage-switch}
\end{equation}
The decision activates one carrier-recovery branch; both branches expose the same complex-sequence interface to the demodulator.

\subsection{Online Statistics and Fixed Parameters}

The sample mean and standard deviation are recomputed from each raw window. The nominal SNR supplies the CV boundary, noise-subtraction term, and reference level for $\hat\gamma_{\mathrm{eff,dB}}$. The CV coefficients, margin $1.10$, proxy floor $10^{-6}$, and 13-dB level are fixed across all runs; the turbulence label, realized gain, recovered data, and branch errors are not control inputs.
